% 自 notes/docs/1151/seminar/W09_1109_預習.md（87808f3）由 tools/md_to_tex.py 於 2026-11-09 轉出一次；此後本檔為正本。
\chapter{導論：三篇的出處與問題}
\label{ch:intro}

\section{三篇的出處與範圍}
\label{sec:intro-1}

\begin{longtable}[]{@{}
  >{\raggedright\arraybackslash}p{(\linewidth - 8\tabcolsep) * \real{0.4100}}
  >{\raggedright\arraybackslash}p{(\linewidth - 8\tabcolsep) * \real{0.2800}}
  >{\raggedright\arraybackslash}p{(\linewidth - 8\tabcolsep) * \real{0.1100}}
  >{\raggedright\arraybackslash}p{(\linewidth - 8\tabcolsep) * \real{0.1000}}
  >{\raggedright\arraybackslash}p{(\linewidth - 8\tabcolsep) * \real{0.1000}}@{}}
\caption{第 9 週的指定文獻。}\tabularnewline
\toprule\noalign{}
\begin{minipage}[b]{\linewidth}\raggedright
篇目
\end{minipage} & \begin{minipage}[b]{\linewidth}\raggedright
原出處
\end{minipage} & \begin{minipage}[b]{\linewidth}\raggedright
課綱指定頁次
\end{minipage} & \begin{minipage}[b]{\linewidth}\raggedright
整理範圍
\end{minipage} & \begin{minipage}[b]{\linewidth}\raggedright
篇幅
\end{minipage} \\
\midrule\noalign{}
\endfirsthead
\toprule\noalign{}
\begin{minipage}[b]{\linewidth}\raggedright
篇目
\end{minipage} & \begin{minipage}[b]{\linewidth}\raggedright
原出處
\end{minipage} & \begin{minipage}[b]{\linewidth}\raggedright
課綱指定頁次
\end{minipage} & \begin{minipage}[b]{\linewidth}\raggedright
整理範圍
\end{minipage} & \begin{minipage}[b]{\linewidth}\raggedright
篇幅
\end{minipage} \\
\midrule\noalign{}
\endhead
\bottomrule\noalign{}
\endlastfoot
Larry Alexander\ix{persons}{alexander@Alexander, Larry} \& Frederick Schauer\ix{persons}{schauer@Schauer, Frederick}, Defending Judicial Supremacy: A Reply & 17 Const. Comment. 455--482 (2000) & 455--473、481--482 & 455--482 全文 & 28 頁，78 註 \\
Larry D. Kramer\ix{persons}{kramer@Kramer, Larry D.}, Judicial Supremacy and the End of Judicial Restraint & 100 Calif. L. Rev. 621--634 (2012) & 621--634 & 621--634 & 14 頁，42 註 \\
John Ferejohn\ix{persons}{ferejohn@Ferejohn, John} \& Pasquale Pasquino\ix{persons}{pasquino@Pasquino, Pasquale}, Constitutional Courts as Deliberative Institutions: Towards an Institutional Theory of Constitutional Justice & in Wojciech Sadurski ed., Constitutional Justice, East and West 21--36 (2002) & 21--36 & 21--36 & 16 頁，22 註 \\
\end{longtable}

Alexander 與 Schauer 一文課綱略去的 474 至 480 頁，是第 II 部分 B 節的後半與第 III 部分的全部，第 2 部照全文整理，各該處標明屬指定頁次以外。

\section{三篇各自處理的問題}
\label{sec:intro-2}

Alexander\ix{persons}{alexander@Alexander, Larry} 與 Schauer\ix{persons}{schauer@Schauer, Frederick} 一文是 1997 年〈On Extrajudicial Constitutional Interpretation〉（110 Harv. L. Rev. 1359）的續篇，回應該文發表後三年間的批評。\footcites{alexander2000}{alexander1997}兩人主張非司法官員應把最高法院的憲法解釋當作與憲法本身同等的權威，理由是法律的核心功能之一在於定紛止爭\ix{terms}{02@\zhlidxcat{定紛止爭論}!定紛止爭}，而憲法文本未能定紛之處，需要一個能以一致的聲音發言的解釋者。第 I 部分論這個問題不能靠憲法文本或制憲史回答，第 II 部分論定紛的價值與最高法院是否為最佳的定紛者，第 III 部分論最高法院的功能是否限於裁判個案。

Kramer\ix{persons}{kramer@Kramer, Larry D.} 一文回應同卷 Posner\ix{persons}{posner@Posner, Richard A.}〈The Rise and Fall of Judicial Self-Restraint〉。\footcites{kramer2012}{posner2012}Posner 把司法自制\ix{terms}{05@\zhlidxcat{解釋理論與自制}!司法自制}的起點歸於 Thayer\ix{persons}{thayer@Thayer, James Bradley} 1893 年的文章\footcite{thayer1893}，Kramer 則追溯到建國時期共和派\ix{terms}{01@\zhlidxcat{司法審查權限之歸屬}!共和派}與聯邦派\ix{terms}{01@\zhlidxcat{司法審查權限之歸屬}!聯邦派}的兩種司法審查\ix{terms}{01@\zhlidxcat{司法審查權限之歸屬}!司法審查}理論，主張兩派都有自制的理由，而各自的理由都以法院以外另有解釋憲法的權威為前提。依 Kramer 的敘述，Brown 案\ix{cases}{02@美國!01@聯邦最高法院!1954Brown v. Board of Education（1954）@Brown v. Board of Education（1954）}之後司法至上\ix{terms}{01@\zhlidxcat{司法審查權限之歸屬}!司法至上}逐漸成為公眾與政治上的共識，自制只剩下保全法院政治資本的審慎\ix{terms}{05@\zhlidxcat{解釋理論與自制}!審慎}考量，左右兩派因此只在自己不在意結果的領域要求法院自制。

Ferejohn\ix{persons}{ferejohn@Ferejohn, John} 與 Pasquino\ix{persons}{pasquino@Pasquino, Pasquale} 比較法國、德國、義大利、西班牙與美國的法院，問的是憲法法院如何說明裁判理由。\footcite{ferejohn2002}兩人區分議會主權\ix{terms}{06@\zhlidxcat{審查制度之類型}!議會主權體制}、權力分立\ix{terms}{06@\zhlidxcat{審查制度之類型}!權力分立體制}與 Kelsen 型\ix{terms}{06@\zhlidxcat{審查制度之類型}!Kelsen 型模式}三種審查制度，認為三者各自產生不同的審議型態：Kelsen\ix{persons}{kelsen@Kelsen, Hans（凱爾森）} 型法院以單一聲音對外說理，內部審議\ix{terms}{07@\zhlidxcat{憲法審判之審議}!內部審議}密集；美國最高法院內部審議少，但以公開的方式處於更廣的外部審議\ix{terms}{07@\zhlidxcat{憲法審判之審議}!外部審議}過程的中心。

\section{與 W08 的銜接}
\label{sec:intro-3}

W08 Tushnet\ix{persons}{tushnet@Tushnet, Mark}〈Non-Judicial Review〉\footcite[註 209]{tushnet2003}的結論以 Alexander\ix{persons}{alexander@Alexander, Larry} 與 Schauer\ix{persons}{schauer@Schauer, Frederick} 為對話對象，註 209 引的是兩人 464 頁「the empirical dimension is one that cannot be avoided」一語所在的段落，Tushnet 據此主張定紛止爭\ix{terms}{02@\zhlidxcat{定紛止爭論}!定紛止爭}論所依賴的經驗前提尚未建立。Alexander 與 Schauer 458 頁註 12 則把 Tushnet 的〈Two Versions of Judicial Supremacy〉\footcite{tushnet1998}列為六篇專門回應之一，472 頁註 48 稱其反對定紛的立場最強；482 頁註 78 承認與批評者交換意見之後，辯論的核心更清楚地呈現為經驗問題，並引 Peabody\ix{persons}{peabody@Peabody, Bruce G.} 說這種研究尚未進行。兩人在 482 頁的結語，以沒有證據顯示其他機構能與最高法院同樣一致地發言為前提，維持最高法院的解釋地位；Tushnet 的結論則以定紛止爭論的經驗前提尚未建立為由，認為非司法審查\ix{terms}{01@\zhlidxcat{司法審查權限之歸屬}!司法審查}的表現不能先被排除。Kramer\ix{persons}{kramer@Kramer, Larry D.} 只在 629 頁註 30 引 Tushnet 的 \en{Taking the Constitution Away from the Courts}\footcite{tushnet1999}，用來支持立法機關內部也有認真的憲法審議；Ferejohn\ix{persons}{ferejohn@Ferejohn, John} 與 Pasquino\ix{persons}{pasquino@Pasquino, Pasquale} 未引 Tushnet。
